% Source: https://artofproblemsolving.com/community/c3h1754490_kinda_stupid_question
% Stable identity: aops-topic-1754490-q01
For a uniform choice on $[0,69]$, the probability that $X$ lies in a subinterval is its length divided by $69$. For every $\varepsilon$ with $0<\varepsilon<1$, the event $X=1$ is contained in the event $1-\varepsilon\le X\le1+\varepsilon$, so
\[
0\le P(X=1)\le\dfrac{2\varepsilon}{69}.
\]
Since this holds for every positive $\varepsilon$, however small, $P(X=1)=\boxed{0}$.

There is no contradiction: probability $0$ does not mean impossible. For a continuous choice every single value has probability $0$ (otherwise, if each of infinitely many points had the same positive probability, the total would exceed $1$); probabilities only become positive for sets of positive length. The expression ``$\dfrac{1}{\infty}$'' is not a number; the precise statement is the limit argument above.
